\textbf{结论名称：}F084\quad 函数与导函数的对称性互推规律

\textbf{为什么要学：}抽象函数没有给出具体式子，却告诉你图像的对称性；借助导数，可以从“左右函数值相等”迅速推到“左右斜率的关系”，也可以反过来判断原函数的对称性。

\textbf{适用条件：}设 $f(x)$ 在 $\mathbb R$ 上可导，$a,b$ 为常数，$t$ 为任意实数。这里的“对称”均指对应图像在整个定义域上满足相应关系。

\textbf{核心结论：}
\[
\boxed{f(a+t)=f(a-t)\ \Longleftrightarrow\ f'(a+t)=-f'(a-t)}
\]
即：$f$ 关于直线 $x=a$ 对称 $\Longleftrightarrow$ $f'$ 关于点 $(a,0)$ 对称。
\[
\boxed{f(a+t)+f(a-t)=2f(a)\ \Longleftrightarrow\ f'(a+t)=f'(a-t)}
\]
即：$f$ 关于点 $(a,f(a))$ 对称 $\Longleftrightarrow$ $f'$ 关于直线 $x=a$ 对称。

\textbf{进一步推广：}若 $f'$ 关于点 $(a,b)$ 对称，即 $f'(a+t)+f'(a-t)=2b$，则
\[
\boxed{f(a+t)-f(a-t)=2bt.}
\]
$b=0$ 时，$f$ 关于 $x=a$ 对称；$b\ne0$ 时，$f$ \textbf{不关于 $x=a$ 对称}，但\textbf{不能断定它没有其他对称轴}。

\textbf{最短记忆：}原函数“同高”$\to$ 导函数“反号”；原函数“围绕中心成对”$\to$ 导函数“同斜率”。

\textbf{例子：}$f(x)=(x-2)^2+1$ 关于 $x=2$ 对称；$f'(x)=2(x-2)$ 关于 $(2,0)$ 对称。

\textbf{考场用法：}先把条件改写成以 $a$ 为中点的 $a+t$ 与 $a-t$；再用上面两组双向等价式，避免猜图像。
